%========================================================================================
%  CV - Maxence TOURNAUD
%  Template    : moderncv (style "banking", couleur "blue")
%  Compilation : pdfLaTeX × 2 passes  |  ou XeLaTeX
%  Distribution: TeX Live / MiKTeX (paquet "moderncv" requis)
%========================================================================================

\documentclass[11pt, a4paper]{moderncv}

% ----------------------------
% STYLE & COULEUR
% ----------------------------
\moderncvstyle{banking}
\moderncvcolor{blue}

% ----------------------------
% RÉDUCTION DE LA TAILLE DE L'EN-TÊTE
% ----------------------------
\renewcommand*{\namefont}{\fontsize{20}{22}\mdseries\upshape}
\renewcommand*{\titlefont}{\fontsize{15}{12}\mdseries\upshape}

% Réduction de l'espace vertical global de l'en-tête
\makeatletter
\renewcommand*{\makecvtitle}{%
  \par\begingroup
    \recomputelengths
    \makecvhead
    \par\vspace{-1.8em}
  \endgroup
}
\makeatother

% ----------------------------
% ENCODAGE & LANGUE
% ----------------------------
\usepackage[utf8]{inputenc}
\usepackage[T1]{fontenc}
\usepackage[french]{babel}

% ----------------------------
% MARGES — réduire les valeurs pour gagner de la place
% ----------------------------
\usepackage[top=1.2cm, bottom=1.2cm, left=1.5cm, right=1.5cm]{geometry}

% ----------------------------
% LARGEUR DE LA COLONNE DES DATES (à gauche)
% ----------------------------
\setlength{\hintscolumnwidth}{2.6cm}

% ----------------------------
% ESPACEMENT GLOBAL PLUS COMPACT
% ----------------------------
% \maketitlewidth n'existe plus dans les versions récentes de moderncv
\AtBeginDocument{\ifdefined\maketitlewidth\setlength{\maketitlewidth}{\textwidth}\fi}

%========================================================================================
% INFORMATIONS PERSONNELLES
%========================================================================================
\name{Maxence}{TOURNAUD}
\title{Étudiant Ingénieur en Alternance Développement Logiciel \& Infrastructure IT}

\address{Orsay (91)}{Île-de-France}{}
\phone[mobile]{07~81~75~69~96}
\email{tournaudmax@gmail.com}
\social[linkedin]{maxence-tournaud}
\social[github]{MaxenceTournaud}

%========================================================================================
\begin{document}
%========================================================================================

% ----------------------------
% EN-TÊTE
% ----------------------------
\makecvtitle

%========================================================================================
% SECTION 1 : PROFIL / RÉSUMÉ PROFESSIONNEL
%========================================================================================
\section{Profil}

\cvitem{}{%
  Étudiant ingénieur à l'\textbf{ENSIIE} en alternance de 3~ans au sein de la DSI du
  \textbf{Paris Brain Institute} (Hôpital Pitié-Salpêtrière, Paris).
  Titulaire d'un \textbf{BUT Informatique} de l'Université Paris-Saclay
  (\textit{Top~10, promotion de 63}), j'ai acquis une expérience concrète en
  \textbf{développement full-stack} (React, Laravel, Angular, .NET) et en
  \textbf{infrastructure IT} (Docker, CI/CD GitLab, Linux, LDAPS).
  Autonome et rigoureux, je m'investis sur des projets à fort impact en
  couvrant l'ensemble du cycle de développement de l'architecture au déploiement en étroite collaboration avec des équipes pluridisciplinaires.%
}

%========================================================================================
% SECTION 2 : EXPÉRIENCES PROFESSIONNELLES
% Ordre chronologique inversé (plus récente en premier)
%========================================================================================
\section{Expériences Professionnelles}

% ----------------------------
% 2.1  ALTERNANCE — Paris Brain Institute / ICM  [EN COURS]
%      Source principale : Rapport d'activité
% ----------------------------
\cventry%
  {Sept.~2025 - Sept.~2028}%
  {Alternant Ingénieur - Développement \& Infrastructure IT}%
  {Paris Brain Institute (ICM) - Direction des Systèmes d'Information}%
  {Paris\,(75), Hôpital Pitié-Salpêtrière}%
  {}% ← Grade / mention (laisser vide ou compléter)
  {%
    \textit{Pôle Infrastructure \& Réseaux (équipes CloudOps et Network).
    Mission principale~: conception, développement et déploiement
    d'une \textbf{CMDB} centralisant l'inventaire du parc informatique
    pour l'ensemble de la DSI (serveurs, réseau, VLANs, logiciels).}%
    \smallskip
    \begin{itemize}
        %-- DÉVELOPPEMENT APPLICATIF --
        \item \textbf{Conception et développement d’une application de CMDB~:}
        Conception et développement d’une application de \textbf{CMDB}, Réalisation de la modélisation des données, de la conception de l’API REST ainsi que de l’interface utilisateur.
    \item \textbf{Sécurité \& gestion des accès~:}
    Conception et implémentation d’un système de \textbf{rôles et permissions}
    adossé à l’\textbf{Active Directory (LDAPS/TLS)}. Mise en place d’un contrôle
    d’accès en lecture/écriture assurant la confidentialité des données sensibles.
    Configuration de \textbf{TLS/LDAPS} pour sécuriser les échanges inter-serveurs.
      %-- DÉPLOIEMENT & DEVOPS --
      \item \textbf{DevOps \& Déploiement automatisé~:}
            Architecturé une infrastructure de \textbf{3~serveurs} sous
            \textbf{Docker} (front-end, back-end, BD) et déployé une chaîne
            \textbf{CI\,/\,CD complète sur GitLab}
    \end{itemize}
  }

% ----------------------------
% 2.2  Stage WebGames (2025 — 4 mois)
% ----------------------------
\cventry%
  {2025\,(4\,mois)}%
  {Stagiaire Développeur JavaScript}%
  {WebGames}%
  {Caen}%
  {}%
  {%
    \begin{itemize}
      \item Développé un \textbf{éditeur de cartes avancé} de plus de
            \textbf{5\,200\,lignes\,JS}~: calques hiérarchiques,
            groupes imbriqués, drag-\&-drop, rotation, sérialisation
            et \textit{grid culling} pour scènes massives -
            \textbf{livré et déployé en production}.
      \item Résolu des défis algorithmiques complexes
            (géométrie, structures de données, culling, responsivité)
            avec livraisons continues.
    \end{itemize}
  }

% ----------------------------
% 2.3  Stage SUEZ Smart Solutions (2024 — 2,5 mois)
% ----------------------------
\cventry%
  {2024\,(2,5\,mois)}%
  {Stagiaire Développeur Full-Stack}%
  {SUEZ Smart Solutions}%
  {Île-de-France}%
  {}%
  {%
    \begin{itemize}
      \item Participé au développement 
            d'\textbf{APIs sécurisées} (cryptographie) pour la gestion
            du traitement d'eau.
      \item Utilisé \textbf{Azure DevOps} en méthodologie agile
            (gestion de projet, pull requests, résolution de conflits).
    \end{itemize}
  }

% ----------------------------
% 2.4  Projets personnels & universitaires (2022-2025)
%      → Commenter ce bloc pour raccourcir le CV à 1 page
% ----------------------------
\cventry%
  {2022 - 2025}%
  {Projets Personnels \& Universitaires}%
  {}{}{}%
  {%
    \begin{itemize}
      \item \textbf{Jeu Battle Royale (Unity\,/\,C\#)} en équipe de 10~:
            sérialisation des inputs joueur, conception du système d'armes ;
            \textbf{responsable du dépôt GitHub}
            (versioning et structuration du code).
        \item \textbf{Home servers Linux~:}
        Déploiement et administration de services auto-hébergés sous Linux, incluant
        un \textbf{serveur web}, ainsi que
        la configuration de \textbf{serveurs Steam} , avec
        automatisation des tâches et du cycle de vie des services grâce à
        \textbf{systemd} et \textbf{cron}.
    \end{itemize}
  }

%========================================================================================
% SECTION 3 : FORMATION
%========================================================================================
\section{Formation}

% ----------------------------
% 3.1  ENSIIE — Diplôme d'ingénieur (en cours)
% ----------------------------
\cventry%
  {Sept.~2025 - Sept.~2028}%
  {Diplôme d'Ingénieur - Informatique \& Mathématiques}%
  {ENSIIE (École Nationale Supérieure d'Informatique
   pour l'Industrie et l'Entreprise)}%
  {Évry-Courcouronnes\,(91)}%
  {}%
  {Formation en alternance, 3\,ans.}

% ----------------------------
% 3.2  BUT Informatique — Université Paris-Saclay (obtenu)
% ----------------------------
\cventry%
  {2022 - 2025}%
  {BUT Informatique (Bac\,+3)}%
  {Université Paris-Saclay, IUT d'Orsay}%
  {Orsay\,(91)}%
  {\textbf{Classement général~: Top\,10} sur une promotion de 63 étudiants}%
  {}

%========================================================================================
% SECTION 5 : LANGUES
%========================================================================================
\section{Langues et Centres d'intérêt}

\cvitemwithcomment{Français}{Langue maternelle}{}
\cvitemwithcomment{Anglais}{Avancé (B2\,/\,C1)}{%
  Lecture technique, communication professionnelle}
  \textbf{Ski alpin de compétition} -
  Titulaire de la \textit{Flèche de Vermeil} (récompense nationale).


%========================================================================================
\end{document}
%========================================================================================